\documentclass[aspectratio=169]{beamer}

% --------- PACKAGES ---------
\usepackage[utf8]{inputenc}
\usepackage[T1]{fontenc}
\usepackage[french]{babel}
\usepackage{graphicx}
\usepackage{booktabs}
\usepackage{hyperref}
\usepackage{tikz}
\usepackage{amsmath}
\usetikzlibrary{arrows.meta, positioning, shapes.geometric}

% --- THEME AND COLOR SCHEME ---
\usetheme{Madrid}
\usecolortheme{beaver}
\definecolor{coverblue}{HTML}{004080}

% --- BEAMER COLOR SETTINGS ---
\setbeamercolor{palette primary}{bg=coverblue,fg=white}
\setbeamercolor{palette secondary}{bg=coverblue,fg=white}
\setbeamercolor{palette tertiary}{bg=coverblue,fg=white}
\setbeamercolor{palette quaternary}{bg=coverblue,fg=white}
\setbeamercolor{structure}{fg=coverblue}
\setbeamercolor{title}{bg=coverblue,fg=white}
\setbeamercolor{frametitle}{bg=coverblue,fg=white}
\setbeamercolor{block title}{bg=coverblue,fg=white}
\setbeamertemplate{navigation symbols}{}
\setbeameroption{hide notes}

% --- TITLE INFORMATION ---
\title[Fiabilite Cluster H100]{Analyse Stochastique de la Fiabilite d'un Cluster NVIDIA H100}
\author[Aarab Ayoub]{Aarab Ayoub \and Equipe Projet}
\institute[UAE ENS Tetouan]{
    ABDELMALEK ESSAADI UNIVERSITY \\
    ENS Tetouan -- Master Data Science \\
    \vspace{0.3cm}
    \textbf{Encadrant:} Pr. Abdellatif El Afia
}
\date{Soutenance 2026}

\begin{document}

% =========================
% SLIDE 1
% =========================
\begin{frame}{Slide 1 -- Title}
    \titlepage
    \vspace{-0.2cm}
    \begin{block}{Short description}
        Objectif: modeliser, estimer et optimiser la fiabilite d'un cluster GPU H100 avec un cadre probabiliste unifie.
    \end{block}
    \begin{alertblock}{Resultat cle}
        KPI principal observe: \textbf{MTTF \(\approx 78.93\) pas}.
    \end{alertblock}
\end{frame}
\note{Je commence par le probleme global et le resultat le plus parlant: le temps moyen avant panne permanente.}

% =========================
% SLIDE 2
% =========================
\begin{frame}{Slide 2 -- Problem \& Motivation}
\begin{columns}[T]
\begin{column}{0.55\textwidth}
\begin{itemize}
    \item Pannes frequentes en entrainement LLM a grande echelle.
    \item Etat reel du GPU non observable directement.
    \item On observe temperature, ECC, XID, puissance.
\end{itemize}
\textbf{Interpretation:} Sans modele probabiliste, la maintenance devient reactive et couteuse.
\end{column}
\begin{column}{0.45\textwidth}
\centering
\begin{tikzpicture}[node distance=0.75cm, scale=0.8, transform shape]
\tikzstyle{box}=[rectangle, rounded corners, draw=coverblue, very thick, align=center, minimum width=3.2cm, minimum height=0.65cm]
\node[box] (h) {Hidden state};
\node[box, below=of h] (o) {Observations};
\node[box, below=of o] (d) {Decision};
\draw[-{Latex[length=2mm]}] (h)--(o);
\draw[-{Latex[length=2mm]}] (o)--(d);
\end{tikzpicture}
\end{column}
\end{columns}
\end{frame}
\note{Je pose la motivation metier: etat cache, observations bruites, et besoin de decision robuste.}

% =========================
% SLIDE 3
% =========================
\begin{frame}{Slide 3 -- DTMC Model}
\begin{block}{Models discrete evolution of system states}
Le DTMC decrit les transitions pas-a-pas entre Healthy, Degraded, MaintenanceRequired, FailedRecoverable et FailedPermanent.
\end{block}
\[
P=\begin{pmatrix}
0.90&0.06&0.02&0.01&0.01\\
0.15&0.70&0.10&0.05&0.00\\
0.40&0.00&0.40&0.15&0.05\\
0.35&0.00&0.20&0.40&0.05\\
0.00&0.00&0.00&0.00&1.00
\end{pmatrix}
\]
\textbf{Interpretation:} Used to estimate failure dynamics de facon simple et interpretable.
\end{frame}
\note{Je presente la matrice comme la loi de dynamique de base qui alimente aussi les autres modeles.}

% =========================
% SLIDE 4
% =========================
\begin{frame}{Slide 4 -- DTMC Result (IMPORTANT)}
\begin{block}{MTTF result}
\[
\mathrm{MTTF}_{Healthy\rightarrow FailedPermanent} \approx 78.93\ \text{pas}
\]
\end{block}
\begin{itemize}
    \item Mesure calculee depuis la matrice de transition executable.
    \item C'est le temps moyen avant panne permanente.
\end{itemize}
\textbf{Interpretation:} Le systeme degrade progressivement avant defaillance finale.
\end{frame}
\note{Je donne une lecture concrete: ce chiffre resume le risque a long terme pour un GPU initialement sain.}

% =========================
% SLIDE 5
% =========================
\begin{frame}{Slide 5 -- CTMC Model}
\begin{block}{Continuous-time modeling using transition rates}
Le CTMC remplace les probabilites par des taux instantanes dans une matrice generatrice \(A\).
\end{block}
\[
A=\begin{pmatrix}
-0.10&0.06&0.02&0.01&0.01\\
0.15&-0.30&0.10&0.05&0.00\\
0.40&0.00&-0.60&0.15&0.05\\
0.35&0.00&0.20&-0.60&0.05\\
0.00&0.00&0.00&0.00&0.00
\end{pmatrix}
\]
\textbf{Interpretation:} On modelise le temps reel, pas seulement des etapes discretes.
\end{frame}
\note{Je precise que CTMC est utile quand les evenements peuvent arriver a n'importe quel instant.}

% =========================
% SLIDE 6
% =========================
\begin{frame}{Slide 6 -- CTMC Result}
\begin{columns}[T]
\begin{column}{0.52\textwidth}
\begin{itemize}
    \item La probabilite de rester operationnel baisse avec le temps.
    \item La courbe est issue de \(P(t)=e^{At}\).
\end{itemize}
\textbf{Interpretation:} Le risque cumule augmente meme sans choc brutal visible.
\end{column}
\begin{column}{0.48\textwidth}
\centering
\includegraphics[width=\textwidth]{results/mandatory_plots/ct_reliability.png}
\end{column}
\end{columns}
\end{frame}
\note{Je lis la courbe en termes de survie: plus on avance en temps, plus la panne permanente devient probable.}

% =========================
% SLIDE 7
% =========================
\begin{frame}{Slide 7 -- Stationary Distribution}
\begin{block}{Long-term behavior of system}
\[
\pi \approx (0.722,\ 0.144,\ 0.048,\ 0.054,\ 0.032)
\]
\[
\text{Uptime long terme} \approx \pi_{Healthy}+\pi_{Degraded} \approx 0.866
\]
\end{block}
\begin{itemize}
    \item La distribution stationnaire resume le regime de long terme.
\end{itemize}
\textbf{Interpretation:} Le systeme reste operationnel la plupart du temps (environ 86\%).
\end{frame}
\note{Je donne l'idee: la stationnaire est un portrait global de fonctionnement a long horizon.}

% =========================
% SLIDE 8
% =========================
\begin{frame}{Slide 8 -- HMM Concept}
\begin{columns}[T]
\begin{column}{0.56\textwidth}
\begin{itemize}
    \item Etat latent non directement observable.
    \item Observations: temperature, ECC, XID, puissance.
    \item HMM relie etat cache et mesures.
\end{itemize}
\textbf{Interpretation:} On transforme des signaux bruites en estimation de sante.
\end{column}
\begin{column}{0.44\textwidth}
\centering
\begin{tikzpicture}[node distance=0.9cm, scale=0.85, transform shape]
\tikzstyle{c}=[circle, draw=coverblue, very thick, minimum size=8mm]
\tikzstyle{o}=[rectangle, draw=coverblue, very thick, rounded corners, minimum width=1.8cm, minimum height=0.55cm]
\node[c] (s1) {$S_{t-1}$};
\node[c, right=of s1] (s2) {$S_t$};
\node[o, below=of s2] (x) {$O_t$};
\draw[-{Latex[length=2mm]}] (s1)--(s2);
\draw[-{Latex[length=2mm]}] (s2)--(x);
\end{tikzpicture}
\end{column}
\end{columns}
\end{frame}
\note{Je montre que l'observation n'est qu'un indice de l'etat reel, pas l'etat lui-meme.}

% =========================
% SLIDE 9
% =========================
\begin{frame}{Slide 9 -- Forward Algorithm}
\begin{block}{Forward = probability / filtering}
\[
\alpha_t(j)=\left(\sum_i \alpha_{t-1}(i)a_{ij}\right)b_j(o_t)
\]
\[
P(S_t=j\mid O_{1:t})=\frac{\alpha_t(j)}{\sum_k\alpha_t(k)}
\]
\end{block}
\begin{itemize}
    \item Calcule la croyance sur chaque etat a chaque instant.
\end{itemize}
\textbf{Interpretation:} Le Forward donne une probabilite, pas une classe unique.
\end{frame}
\note{Je souligne que le Forward sert au suivi en ligne et a la quantification d'incertitude.}

% =========================
% SLIDE 10
% =========================
\begin{frame}{Slide 10 -- Viterbi Algorithm}
\begin{block}{Best path reconstruction}
\[
\hat{s}_{1:T}=\arg\max_{s_{1:T}} P(s_{1:T},o_{1:T})
\]
\end{block}
\begin{itemize}
    \item True state = ground truth simule.
    \item Viterbi state = sequence inferee.
\end{itemize}
\textbf{Interpretation:} Viterbi reconstruit la trajectoire la plus plausible des etats caches.
\end{frame}
\note{Je distingue bien filtrage (Forward) et reconstruction globale (Viterbi).}

% =========================
% SLIDE 11
% =========================
\begin{frame}{Slide 11 -- HMM Result}
\begin{columns}[T]
\begin{column}{0.52\textwidth}
\begin{itemize}
    \item Posterior Forward suit les changements de regime.
    \item true\_state et viterbi\_state sont proches dans les zones stables.
\end{itemize}
\textbf{Interpretation:} Le HMM reconstruit correctement les regimes latents dominants.
\end{column}
\begin{column}{0.48\textwidth}
\centering
\includegraphics[width=\textwidth]{results/mandatory_plots/hmm_forward_posterior.png}
\end{column}
\end{columns}
\end{frame}
\note{Je lis les pics de probabilite comme des moments de forte certitude diagnostique.}

% =========================
% SLIDE 12
% =========================
\begin{frame}{Slide 12 -- MDP Model}
\begin{block}{Decision-making under uncertainty}
\[
V^*(s)=\max_a\left[R(s,a)+\gamma\sum_{s'}P(s'\mid s,a)V^*(s')\right]
\]
\end{block}
\begin{itemize}
    \item Le MDP choisit l'action qui maximise valeur immediate + valeur future.
\end{itemize}
\textbf{Interpretation:} Politique optimale = regle de decision par etat de sante.
\end{frame}
\note{Je presente Bellman comme une balance entre present et futur sous incertitude.}

% =========================
% SLIDE 13
% =========================
\begin{frame}{Slide 13 -- MDP Result (IMPORTANT)}
\begin{block}{Policy table: state \(\rightarrow\) action}
\begin{tabular}{ll}
\toprule
Etat & Action optimale \\
\midrule
Healthy & keep\_running \\
Degraded & power\_cap \\
MaintenanceRequired & isolate\_maintain \\
FailedRecoverable & reboot \\
FailedPermanent & isolate\_maintain \\
\bottomrule
\end{tabular}
\end{block}
\textbf{Interpretation:} Chaque etat a une action cible claire et operationnelle.
\end{frame}
\note{Je relie chaque action a une logique metier concrete de fiabilite.}

% =========================
% SLIDE 14
% =========================
\begin{frame}{Slide 14 -- Reinforcement Learning}
\begin{block}{Learning without knowing the model exactly}
\[
Q(s,a)\leftarrow Q(s,a)+\alpha\left[r+\gamma\max_{a'}Q(s',a')-Q(s,a)\right]
\]
\end{block}
\begin{itemize}
    \item Q-learning et SARSA apprennent par interaction.
    \item Pas besoin d'un modele parfait en entree.
\end{itemize}
\textbf{Interpretation:} RL complete le MDP quand le monde reel est partiellement connu.
\end{frame}
\note{Je rappelle que RL est utile quand les dynamiques reelles evoluent ou sont incompletes.}

% =========================
% SLIDE 15
% =========================
\begin{frame}{Slide 15 -- RL Results}
\begin{columns}[T]
\begin{column}{0.54\textwidth}
\begin{tabular}{lc}
\toprule
Methode & Retour moyen \\
\midrule
Q-learning & 63.87 \\
SARSA & 67.21 \\
\bottomrule
\end{tabular}

\vspace{0.2cm}
\textbf{Interpretation:} L'agent apprend un comportement performant; SARSA est legerement meilleur ici.
\end{column}
\begin{column}{0.46\textwidth}
\centering
\includegraphics[width=\textwidth]{results/mandatory_plots/rl_convergence.png}
\end{column}
\end{columns}
\end{frame}
\note{Je commente la tendance globale plutot que les oscillations episode par episode.}

% =========================
% SLIDE 16
% =========================
\begin{frame}{Slide 16 -- POMDP / Non-Markov}
\begin{block}{Handles uncertainty and memory}
\[
b_{t+1}(s')=\eta\,O(o_{t+1}\mid s',a_t)\sum_s T(s'\mid s,a_t)b_t(s)
\]
\[
z_t=(\tau_t,c_t)\quad \text{(age + cumul ECC)}
\]
\end{block}
\begin{itemize}
    \item POMDP: decision sur croyance.
    \item Non-Markov: le risque depend aussi de l'historique.
\end{itemize}
\textbf{Interpretation:} Plus d'incertitude et de memoire = politique plus prudente.
\end{frame}
\note{Je montre que cette extension rapproche le modele des conditions operationnelles reelles.}

% =========================
% SLIDE 17
% =========================
\begin{frame}{Slide 17 -- Final Insights}
\begin{columns}[T]
\begin{column}{0.55\textwidth}
\begin{itemize}
    \item DTMC/CTMC: dynamique de fiabilite.
    \item HMM: estimation d'etat latent.
    \item MDP/RL: decision et apprentissage.
\end{itemize}
\textbf{Interpretation:} Un seul pipeline: \textbf{DTMC \(\rightarrow\) HMM \(\rightarrow\) MDP \(\rightarrow\) RL}.
\end{column}
\begin{column}{0.45\textwidth}
\centering
\begin{tikzpicture}[node distance=0.65cm, scale=0.8, transform shape]
\tikzstyle{b}=[rectangle, rounded corners, draw=coverblue, very thick, align=center, minimum width=2.5cm, minimum height=0.6cm]
\node[b] (d) {DTMC / CTMC};
\node[b, below=of d] (h) {HMM};
\node[b, below=of h] (m) {MDP};
\node[b, below=of m] (r) {RL};
\draw[-{Latex[length=2mm]}] (d)--(h);
\draw[-{Latex[length=2mm]}] (h)--(m);
\draw[-{Latex[length=2mm]}] (m)--(r);
\end{tikzpicture}
\end{column}
\end{columns}
\end{frame}
\note{Je termine l'histoire du projet: chaque bloc repond a une question differente mais complementaire.}

% =========================
% SLIDE 18
% =========================
\begin{frame}{Slide 18 -- Conclusion}
\begin{block}{What we learned}
\begin{itemize}
    \item Fiabilite quantifiee: MTTF \(\approx 78.93\).
    \item Etat latent estimable via HMM.
    \item Decision robuste via MDP + RL.
\end{itemize}
\end{block}
\begin{alertblock}{Why this approach is powerful}
Cadre unique, interpretable et extensible pour passer de la prediction a l'action.
\end{alertblock}
\textbf{Next step:} calibration sur telemetrie reelle DCGM/XID.
\end{frame}
\note{Message final simple: on sait modeliser, estimer et agir; la suite est la validation sur donnees reelles.}

\end{document}
